\documentclass{article} %210 mm × 297 mm
\usepackage[utf8]{inputenc}
\usepackage[a4paper,left=1.5cm, right=1.5cm, top=2cm, bottom=2cm]{geometry}
\usepackage{graphicx}
%\usepackage{blindtext}
\usepackage{multirow}
\usepackage{mhchem}
\usepackage{url}
\usepackage{subfigure}
\usepackage{amsfonts,latexsym} % para tener disponibilidad de diversos simbolos
\usepackage{enumerate}
%\usepackage{booktabs}
\usepackage{float}
%\usepackage{threeparttable} 
\usepackage{array,colortbl}
%\usepackage{ifpdf}
%\usepackage{rotating}
\usepackage{cite}
\usepackage{stfloats}
\usepackage{url}
\usepackage{listings}
\usepackage{fancyhdr}

\usepackage{color}
%\usepackage{hyperref}
%\usepackage{wrapfig}
\usepackage{array}
%\usepackage{multirow}
%\usepackage{adjustbox}
%\usepackage{nccmath}
\usepackage[dvipsnames]{xcolor}


\newcommand{\titulo}{Abgabe 9; Statistik}
\newcommand{\nombre}{Liliana Sanfilippo}
\newcommand{\FCE}{\textbf{SoSe 2025}}
\newcommand{\dat}{\textbf{Abgabe 19.06.2025}}

\newcommand{\ma}[1]{\( #1 \)}
\newcommand{\mab}[1]{\[ #1 \]}
\newcommand{\kl}[1]{\{ #1 \}}
\newcommand{\p}[0]{\mathbb{P}}
\newcommand{\E}[1]{\mathbb{E}[#1]}
\newcommand{\sumEN}[0]{\sum_{i = 1}^n}

\usepackage{fancyhdr}
\pagestyle{fancy}
\fancyhead{}
\fancyfoot{}
\fancyhead[RO]{\small\dat}
\fancyhead[LO]{\small\FCE}
\fancyfoot[RO,LE] {\thepage}

\setcounter{page}{1} %Número de la página, la editorial lo cambiará



\usepackage[onehalfspacing]{setspace}
\begin{document}


\begin{tabular}{p{12cm}}
{{\Large\bf\titulo}}\\
\vspace{0.2cm}\\
 {\large\nombre}\\
\end{tabular}
\vspace{0.2cm}\\


\section*{Aufgabe 2}
\subsection*{a)}
\ma{\Omega =\kl{0,1}^n} mit \ma{\omega = 1} bedeutet, dass eine Zahl geworfen wurde. 
\mab{\omega \in \Omega := (\omega_1, ..., \omega_n)}
\mab{\p(\omega_i) = \sumEN \frac{1}{2} \forall i \in \mathbb{N}, \quad \p (\omega) =  \p(\omega_i) =  \frac{1}{2^n}}
\mab{S_n(\omega) = \sumEN \omega_i, \quad S_n \sim \mathcal{B}(n, \frac{1}{2})}
Annahmen:
\begin{itemize}
    \item Jeder Münzwurf ist unabhängig vom vorherigen
    \item evtl. jeder Wurf ist gleich verteilt
\end{itemize}

\subsection*{b)}Hi rnochmal gucken, ob das mit dem $<$ und $\leq$ alles so passt

\mab{\p(-\frac{1}{20} < \frac{1}{n} S_n - \frac{1}{2} \leq \frac{1}{20}) \geq \frac{19}{20} 
\quad \Leftrightarrow \quad 
\p( |\frac{1}{n} S_n  - \frac{1}{2}| \leq \frac{1}{20}) \geq \frac{19}{20}
\quad \Leftrightarrow \quad 
\p(|\frac{S_n}{n} - \E{\frac{S_n}{n}}| \geq \frac{1}{20} ) \leq 1 - \frac{19}{20} }

\mab{
\p(|\frac{S_n}{n} - \frac{1}{2}| \geq \frac{1}{20} ) \leq \frac{Var(\frac{S_n}{n})}{(\frac{1}{20})^2}
= \frac{\frac{1}{4n}}{(\frac{1}{20})^2}
= \frac{\frac{1}{4n}}{400}
= \frac{100}{n}
}

\mab{\frac{100}{n} \leq \frac{1}{20} \Rightarrow n \geq 2000}


\subsection*{c)}
\mab{
\p(-\frac{1}{20} \leq \frac{1}{n} S_n - \frac{1}{2} \leq \frac{1}{20}) 
=
\p(-\frac{1}{20} + \frac{1}{2} \leq \frac{1}{n} S_n \leq \frac{1}{20} + \frac{1}{2}) 
= 
\p(\frac{9}{20} \leq \frac{1}{n} S_n \leq \frac{11}{20})  
= 
\p(\frac{9n}{20} \leq  S_n \leq \frac{11n}{20})
}
\mab{\p(\frac{9n}{20} -1  <  S_n \leq \frac{11n}{20})}
\mab{
\p(a < \frac{S_n - np}{\sqrt{npq}} \leq b) = \p(\frac{\frac{9n}{20} -1  - np}{\sqrt{npq}} < \frac{S_n - np}{\sqrt{npq}} \leq \frac{\frac{11n}{20} - np}{\sqrt{npq}}) \approx  \Phi(\frac{\frac{11n}{20} - np}{\sqrt{npq}}) - \Phi(\frac{\frac{9n}{20} -1  - np}{\sqrt{npq}}) 
}
\mab{ = \Phi(\frac{\frac{11n}{20} - n\frac{1}{2}}{\sqrt{n\frac{1}{2}\frac{1}{2}}}) - \Phi(\frac{\frac{9n}{20} -1  - n\frac{1}{2}}{\sqrt{n\frac{1}{2}\frac{1}{2}}})
= \Phi(\frac{\frac{11n}{20} - n\frac{1}{2}}{\sqrt{n\frac{1}{2}\frac{1}{2}}}) - \Phi(\frac{\frac{9n}{20} -1  - n\frac{1}{2}}{\sqrt{n\frac{1}{2}\frac{1}{2}}})
}
\end{document}